% ============================================================
%  2. Analyse et Conception  (~3 pages)
% ============================================================
\section{Analyse et Conception}

\subsection{Modèle conceptuel}
% Insérer le diagramme UML conceptuel
% Expliquer les choix : Tuple, Attribut, Valeur abstraite

Avant de concevoir la base de données, nous avons défini un diagramme conceptuel qui fait complétement abstraction des toutes implémentations physiques relative à Java.

\begin{figure}[H]
    \centering
    \includegraphics[width=0.8\textwidth]{diagramme_conceptuel.png}
    \caption{Modèle conceptuel de données}
    \label{fig:conceptuel}
\end{figure}

\subsection{Modèle physique}
% Insérer le diagramme UML physique complet
% Justifier les choix d'implémentation

\begin{figure}[H]
    \centering
    \includegraphics[width=\textwidth]{diagramme_physique.png}
    \caption{Modèle physique — diagramme de classes complet}
    \label{fig:physique}
\end{figure}

\subsection{Choix de conception notables}

\begin{itemize}
    \item \textbf{Compteur SERIAL global} : le compteur est centralisé dans
          \texttt{BaseDeDonnees} et partagé entre toutes les tables elles ont accès au \textit{compteurSerial}.
    \item \textbf{Hiérarchie de Valeur} : \texttt{ValeurInt} et \texttt{ValeurVarchar}
          héritent de la classe abstraite \texttt{Valeur}.
    \item \textbf{Immutabilité} : les listes exposées par \texttt{Table} et
          \texttt{Tuple} sont immuables.
          
    \item \textbf{Classe Table} : Nous nous sommes questionnées si on devrait mettre les Expressions logiques dans \textbf{supprimerTuple} mais nous avons utilisé \textit{les tuples plutôt} afin de réduire le couplage entre Expressions logique et Base de Données. %Utilisation des structures et des extends list %
    
\end{itemize}
